% Section 2, study design (restructure plan, part 2). Sources: experiments/protocol.md (§2, §4, §15.4, §15.5, §15.7,
% §18); docs/decision_log.md (the D-ids in the text); scripts/propose_labels.py (d19, combine_repeats; Table
% tab:label-rules); src/tracecontract/audit.py (the audit's checks and blocking findings); the generator template
% experiments/prompts/workflows/tracecontract-generate__wf_b796bb4f-1dd.js (the over-fitting sentence). Every number is a
% macro: \fv (evidence_ledger_fresh.md), \dv (evidence_ledger_design.md), \nrv (evidence_ledger_natural_run.md).
\section{Study design}
\label{sec:design}
This section describes the controlled study. The natural arm follows the same design with the differences that \S\ref{sec:natural-cases} lists, after a screening funnel of deterministic stages S1 to S8a, a model-judged stage S9 and a replay stage S10, whose families S11 accepts for building (Table~\ref{tab:natural-funnel}); each accepted family is then built in phases P0 to P7 (Appendix~\ref{app:natural}). ``Protocol~\S$n$'' refers to the locally logged protocol (\texttt{experiments/protocol.md}), and ``D$n$'' to an entry of the decision log (\texttt{docs/decision\_log.md}).

\subsection{Tasks, sessions and attempts}
\label{sec:design-tasks}
The population is controlled (protocol~\S2). Each task family is a real agent session on a pinned commit of a permissively licensed Python library, whose user turns the study's own agent wrote and hashed before the session ran (protocol CS-v1). Its feedback is therefore synthetic, and we never describe it as user traffic. A scenario fixes $U_1$, $U_2$, a neutral reply to clarification questions, and the intended requirements with the turn that introduced each (Table~\ref{tab:scenarios}). Figure~\ref{fig:design} shows the design.
\begin{itemize}
\item CS-01 (\texttt{humanize}): $U_1$ asks \texttt{naturalsize} to drop a trailing \texttt{.0}; the correction makes this opt-in and keeps the default and custom-format output, retracting part of $U_1$.
\item CS-02 (\texttt{tinydb}): atomic JSON writes; the correction adds read-after-write, unchanged bytes, no leftover temporary files and an intact original on failure.
\item CS-03 (\texttt{inflection}): \texttt{titleize} keeps acronyms; the correction keeps all-uppercase input title-cased and \texttt{humanize} and \texttt{underscore} unchanged.
\item CS-04 (\texttt{more-itertools}): \texttt{consecutive\_groups} gets a \texttt{step}; the correction adds negative steps, an error for \texttt{step=0}, positional ordering, lazy groups and an unchanged default.
\end{itemize}
The session agent, a subagent of an AI coding agent, works in a fresh git copy of the base commit and runs repository code only through wrappers that execute it in an offline sandbox. Its workspace is captured after $U_1$ (the pre-correction state) and at the end (the reference); the correction arrives through a messaging channel that adds a short harness frame (D10). We call the agent session that ran the study, and launched every other agent, the \emph{build session}. The \emph{owner} is the person who owns the study's repository and instructs the build session. A \emph{builder decision} is a choice that the build session made and logged in the decision log without the owner; the paper marks those that await the owner's review (Appendix~\ref{app:deviations}).

Fresh agents, subagents of the same AI coding agent, attempt each task in \fv{run.configurations.w} configurations, A and B (protocol~\S15.4): A is the build session's default model and B the smallest model tier the agent tool offers, \fv{run.configurations.w} model tiers of one model family of one provider. The protocol and the decision log name the agent tool but no model identifier, which the protocol keeps in private run records. The repository does name models elsewhere: each of its commit messages carries a co-author line that names a model, and its committed controlled-session transcripts carry the session agent's model identifier in their raw records (D79). The paper names no model. Configuration B also writes the independent labeler (Table~\ref{tab:roles}).

An attempt receives $U_1$ (\texttt{u1}), or $U_1$, the line ``Later in the session the user added:'' and $U_2$ (\texttt{final}), each verbatim up to recorded path substitutions (Appendix~\ref{app:deviations}), and nothing else: no reference, transcript, verifier, candidate or label (protocol~\S15.4). Each attempt works in its own workspace (a fresh git repository of the base snapshot, the instruction and the sessions' sandbox wrappers) for at most \dv{design.attempt.minutes} minutes. With $k=\dv{design.k}$ per family, variant and configuration there are \fv{attempts} attempts. An attempt's final state, without the files the patch policy removes (its test-file edits among them), is its candidate. Each transcript is audited for access outside the workspace. Under the review rule of protocol v0.10, fixed after the first family's flags had been read and before any grading, \fv{attempts.excluded.w} attempt is excluded: \texttt{\fv{excluded.1}} wrote, read and ran helper files of its own in the build session's scratch directory (D62). A sensitivity view leaves out all \fv{sens.excluded} flagged attempts (protocol~\S18.1). The attempts' context also held harness material that holds no reference, transcript, verifier, candidate or label: an environment block naming the study repository and that scratch directory, and task-list reminders with the titles (and, in CS-03 and CS-04, the descriptions) of the build session's task-list items, in the attempts whose audit records them (protocol~\S18.2; Appendix~\ref{app:deviations}).

\subsection{Verifier conditions}
\label{sec:design-conditions}
The same proposer prompt (\texttt{contract:v1}) recovers the correction-aware contract from the whole session and the initial-only contract from $U_1$ alone. Superseded and unresolved requirements are never scored. A generator writes black-box tests from a contract (the prompt asks for \dv{design.tests.min} to \dv{design.tests.max}; C1, C3 and B2 have \fv{checks.plain.min} to \fv{checks.plain.max} scored checks, Table~\ref{tab:verifiers}): C1 from the initial-only and C3 from the correction-aware contract, both single-pass. B2's generator gets the whole transcript and no contract. Every generator gets the same template (apart from the deviations of Appendix~\ref{app:deviations}), the base repository and the reference diff, with the instruction ``Do not over-fit to its particular implementation choices''. The C1--C3 contrast therefore isolates what the contract asks the generator to test, not what information reached it. B1 is the repository's tests at the reference state that pass on the reference; since the reference is the session agent's final state, they include tests that agent added. In the natural arm only, a secondary arm B3 gets the initial request and the reference diff and no contract, and keeps a test only if it fails on the base and passes on the reference (protocol~\S15.6).

C2 and C4 start from C1's and C3's tests and apply the audit (\texttt{audit:v1}), which grades audit candidates only (the development candidates, never a held-out candidate): the reference \dv{design.audit.reference.repeats.w} times, the no-op (the base snapshot unchanged), \dv{design.audit.attacks.w} attacks written from the same contract, each to violate a requirement, and \dv{design.audit.alternatives.w} valid alternative. The audit passes when the reference passes every repeat, the no-op and every attack fail, the alternative passes, every scored requirement maps to a check and a static lint passes. It also asks each attack to fail a check mapped to the requirement it targets, and reports an attack that fails only other checks as a finding that does not fail the audit. After a failing audit a reviser changes the tests, keeping \dv{design.revise.tests.min} to \dv{design.revise.tests.max} of them, at most \dv{design.audit.revisions.w} times (C2 and C4 have \fv{checks.audited.min} to \fv{checks.audited.max} scored checks), and a verifier whose audit still fails is frozen as it stands (D27). The audit's attacks come from the verifier's own contract, so it cannot supply a requirement the contract lacks, though a revision can tighten checks on behaviour the contract leaves open, as CS-03's did (\S\ref{sec:results-nulls}). Every verifier that grades the attempts, the label-support verifiers included, was frozen by content hash before the first attempt (Table~\ref{tab:roles}).

\begin{table}[t]\centering\scriptsize\setlength{\tabcolsep}{3pt}
\caption{Who wrote what. Configuration: A or B where a record holds it (A and B use different models of one model family of one provider, which are not named; protocol~\S15.4, \S15.5, \S19.2; D64); not named: a record holds the identifier, and the paper does not name it; not recorded: no record holds it, so the statement that these roles use the same model family cannot be checked from the release. Frozen before: what did not exist yet when the output was frozen.}\label{tab:roles}
\begin{tabular}{@{}>{\raggedright\arraybackslash}p{0.2\linewidth}>{\raggedright\arraybackslash}p{0.13\linewidth}>{\raggedright\arraybackslash}p{0.24\linewidth}>{\raggedright\arraybackslash}p{0.22\linewidth}>{\raggedright\arraybackslash}p{0.15\linewidth}@{}}\toprule
Role & Configuration & Inputs & Outputs & Frozen before \\\midrule
\emph{Controlled arm} & & & & \\
Scenario writer (the study's agent) & same model as the session agent & the study's design & $U_1$, $U_2$, intended requirements; the reference contract's quotes & its session (hashed) \\
Session agent & not named & $U_1$, then $U_2$; the base repository & pre-correction state, reference, transcript & --- \\
Contract proposers & not recorded & $U_1$ alone, or the whole session & initial-only and correction-aware contracts & the verifiers \\
Generators and revisers of C1--C4 and B2 & not recorded & a contract (B2: the transcript), the base repository, the reference diff; revisers: the audit's findings & tests & any held-out candidate \\
Reference-check writers & A & the reference contract, the base repository, the reference diff & reference checks & the attempts \\
Independent labeler & B & the reference contract, the base snapshot & per-requirement checks & the attempts \\
Fresh attempts & A and B & $U_1$, or $U_1$ and $U_2$; the base snapshot & attempts & --- \\
Pilot pool writers and mappers & not recorded & the reference instruction or the reference; the contracts & constructed candidates; coverage maps & --- \\\addlinespace
\emph{Natural arm} & & & & \\
S9 judges & A & a pseudonymised comment, issue and diffs & screening verdicts & --- \\
Proposers, generators, revisers, reference-contract drafter, reference-check writers and reviewers, pool writers, mappers & A & the issue or the thread, the diffs, the base repository & contracts, verifiers, reference contract and checks, held-out pool & the next phase (D64) \\
Independent labeler and its revision & B & the rendered reference contract, the base snapshot & per-requirement checks & the held-out pool \\
\bottomrule\end{tabular}\end{table}

\subsection{Labels}
\label{sec:design-labels}
The labeling standard is the reference contract plus a regression norm. A controlled reference contract states the study's intended requirements, hashed before each session ran; the study's agent operationalised it with quotes from the session, and only a model reviewed it. The label-support verifiers, which are not conditions, implement the standard.
\begin{itemize}
\item The \emph{reference checks} were written by configuration-A agents from the reference contract and validated only on the reference and the no-op. Only CS-02's went through a review-and-fix loop (D33).
\item The \emph{regression norm} is the base commit's tests that pass on the reference; a candidate that fails any of them is invalid.
\item The \emph{independent labeler} is written per family by a configuration-B agent from the reference contract and the base snapshot alone, with no reference diff, reference check, verifier, candidate or label, and no regression part (protocol~\S15.5). It gives a second labeling, reported beside the label rule's and never changing it; its agreement with the label rule is agreement between models that read the same reference contract. Like the reference checks, it was validated on the reference, which it had to pass (protocol~\S15.5) and passes in every repeat (Appendix~\ref{app:construction}), and the regression norm is the base tests that pass on the reference; so, as in the natural arm, every label source treats the session agent's final state as correct on what it tests.
\end{itemize}
The label rule labels an attempt invalid if a reference check or the regression norm fails, and valid if both pass (Table~\ref{tab:label-rules}). A violation the reference checks do not exercise is therefore labelled valid for attempts, so no verifier is charged with accepting it. The labels are not independent of the verifiers that score best: by the study's procedure, the reference checks and the correction-aware verifiers C3 and C4 are written by models of the same model family, from the same reference diff, and from requirement lists that the mappers judge to cover the same requirements (the correction-aware contracts cover every requirement of the reference contracts, Table~\ref{tab:coverage}), so their blind spots are likely correlated (Appendix~\ref{app:pilot}). That C3 and C4 accept none of the correction violators (\fv{cv.c3}, \fv{cv.c4}) is therefore partly agreement between renderings of the same requirements. No human label exists; the blinded review packet exists and has not been used.

\begin{table}[t]\centering\scriptsize
\caption{Label rules by candidate source (rule D19, \texttt{scripts/propose\_labels.py}; protocol~\S15.5, \S15.6; D64). D19 combines the reference checks, the regression norm and, for constructed candidates, the construction intent recorded when the candidate was built. An ambiguous candidate is left out of FA and FR and queued for human review.}\label{tab:label-rules}
\begin{tabular}{@{}>{\raggedright\arraybackslash}p{0.27\linewidth}>{\raggedright\arraybackslash}p{0.69\linewidth}@{}}\toprule
Candidates & Rule \\\midrule
Constructed: the pilot's alternatives and near misses; the natural pool's alternatives and near misses & Valid only if built valid and both the reference checks and the norm pass, each with a decision. Invalid only if built invalid and some reference check fails, and either a failing check maps to a targeted requirement or REGRESSION is targeted and the norm fails. Anything else is ambiguous. \\
Observed: the pilot's pre-correction states; fresh attempts & Invalid if a reference check or the norm fails; valid if both pass; ambiguous only on a missing decision. \\
Natural pre-review states & Invalid only through a failing reference check; ambiguous if the reference checks pass, whatever the norm says; never valid. \\
Every natural candidate & The reference checks and the norm are graded \dv{natural.d19.repeats.w} times; repeats that disagree make the label ambiguous, and unusable reference checks make every D19 label of the family ambiguous. \\
Second labeling (the independent labeler): attempts and natural candidates & Reported beside D19, never changing it. A labeler pass on a natural pre-review state is ambiguous; an unusable labeler is left out (NP-03, NP-19). \\
Source-evidence label: natural pre-review states & Invalid with respect to the correction, on the strength of the reviewer's request, the addressing commit and the agreement of both model judges of S9. \\
\bottomrule\end{tabular}\end{table}

\subsection{Metrics and plan}
\label{sec:design-metrics}
False acceptance is the share of invalid held-out candidates whose outcome is a pass, false rejection the share of valid ones without a pass; every other outcome counts as a rejection. Each distinct graded patch counts once: \fv{units} graded patches come from the \fv{attempts} attempts, less the excluded attempt. The unit of analysis is the task family. The \fv{cmp.count} planned paired comparisons, on FA and on FR, are C1--C3 and C2--C4 (initial-only against correction-aware), C1--C2 and C3--C4 (the audit), B2--C3 and B2--C4 (structured contract against the session-LLM baseline) and B1--C4 (repository tests). Each has an exact sign-flip p-value over families: with \fv{families.w} families there are \dv{design.signflip.assignments} sign assignments, so no p-value can fall below \fv{cmp.floor.one} one-sided or \fv{cmp.floor.two} two-sided, and a family with no difference raises the floor.

The per-family table is the primary display. The reference-gated view (the protocol's oracle-gated view), which drops a verifier in each family where it fails its own reference and reports the lost yield, is secondary; both were planned before grading (protocol v0.7, \S15.7). The silent-loss subset of \S\ref{sec:results-controlled} is a post hoc reading of that secondary view, chosen after the results and an internal review. The analysis plan is logged locally: hashed before the first evaluation run and amended \dv{protocol.amendments} times, each version hashed and timestamped (Table~\ref{tab:amendments}). \texttt{\FRESHRUN} tests, on new candidates, a design tuned on the same \fv{families.w} families.
